\documentclass[10pt,a4paper]{amsart}
\usepackage{amsmath,amssymb,amsthm,graphicx}
\usepackage{verbatim}
\newenvironment{reference}{\comment}{\endcomment}
%\newenvironment{reference}{}{}
\newenvironment{slogan}{\comment}{\endcomment}
\newenvironment{history}{\comment}{\endcomment}

% For commutative diagrams we use Xy-pic
\usepackage[all]{xy}

% We use 2cell for 2-commutative diagrams.
\xyoption{2cell}
% Theorem environments.
%
\theoremstyle{plain}
\newtheorem{theorem}[subsection]{Theorem}
\newtheorem{proposition}[subsection]{Proposition}
\newtheorem{lemma}[subsection]{Lemma}

\theoremstyle{definition}
\newtheorem{definition}[subsection]{Definition}
\newtheorem{example}[subsection]{Example}
\newtheorem{exercise}[subsection]{Exercise}
\newtheorem{situation}[subsection]{Situation}

\theoremstyle{remark}
\newtheorem{remark}[subsection]{Remark}
\newtheorem{remarks}[subsection]{Remarks}

\numberwithin{equation}{subsection}

% Macros
%
\def\lim{\mathop{\mathrm{lim}}\nolimits}
\def\colim{\mathop{\mathrm{colim}}\nolimits}
\def\Spec{\mathop{\mathrm{Spec}}}
\def\Hom{\mathop{\mathrm{Hom}}\nolimits}
\def\Ext{\mathop{\mathrm{Ext}}\nolimits}
\def\SheafHom{\mathop{\mathcal{H}\!\mathit{om}}\nolimits}
\def\SheafExt{\mathop{\mathcal{E}\!\mathit{xt}}\nolimits}
\def\Sch{\mathit{Sch}}
\def\Mor{\mathop{\mathrm{Mor}}\nolimits}
\def\Ob{\mathop{\mathrm{Ob}}\nolimits}
\def\Sh{\mathop{\mathit{Sh}}\nolimits}
\def\NL{\mathop{N\!L}\nolimits}
\def\CH{\mathop{\mathrm{CH}}\nolimits}
\def\proetale{{pro\text{-}\acute{e}tale}}
\def\etale{{\acute{e}tale}}
\def\QCoh{\mathit{QCoh}}
\def\Ker{\mathop{\mathrm{Ker}}}
\def\Im{\mathop{\mathrm{Im}}}
\def\Coker{\mathop{\mathrm{Coker}}}
\def\Coim{\mathop{\mathrm{Coim}}}

% Boxtimes
%
\DeclareMathSymbol{\boxtimes}{\mathbin}{AMSa}{"02}

%
% Macros for moduli stacks/spaces
%
\def\QCohstack{\mathcal{QC}\!\mathit{oh}}
\def\Cohstack{\mathcal{C}\!\mathit{oh}}
\def\Spacesstack{\mathcal{S}\!\mathit{paces}}
\def\Quotfunctor{\mathrm{Quot}}
\def\Hilbfunctor{\mathrm{Hilb}}
\def\Curvesstack{\mathcal{C}\!\mathit{urves}}
\def\Polarizedstack{\mathcal{P}\!\mathit{olarized}}
\def\Complexesstack{\mathcal{C}\!\mathit{omplexes}}
% \Pic is the operator that assigns to X its picard group, usage \Pic(X)
% \Picardstack_{X/B} denotes the Picard stack of X over B
% \Picardfunctor_{X/B} denotes the Picard functor of X over B
\def\Pic{\mathop{\mathrm{Pic}}\nolimits}
\def\Picardstack{\mathcal{P}\!\mathit{ic}}
\def\Picardfunctor{\mathrm{Pic}}
\def\Deformationcategory{\mathcal{D}\!\mathit{ef}}
\usepackage{cleveref}
\crefname{originallemma}{original lemma}{original lemmas}
\Crefname{originallemma}{Original lemma}{Original lemmas}
\makeatletter
\newlabel{lemma-constant-exact}{{093J}{source}{}{sourcepointer-1}{}}
\newlabel{lemma-constant-exact@cref}{{[originallemma][0][]093J}{[1][0][]source}{}{}{}}
\makeatother
\title{Sections of a constant tensor product: \textbackslash{}underline\{M \textbackslash{}otimes\_\textbackslash{}Lambda Q\}(U) = \textbackslash{}underline\{M\}(U) \textbackslash{}otimes\_\textbackslash{}Lambda Q for finitely presented Q}
\author{The Stacks Project Authors}
\date{stacks-project:sites-modules.tex:local-snapshot-20261012; GNU FDL 1.2 or later}
\begin{document}
\maketitle
\noindent\textit{Editorial note.} All delivered mathematical text is the author's own, copied verbatim from the pinned Stacks Project chapter file sites-modules.tex; no statement or proof step was rewritten, reordered or removed.\par\medskip
\noindent\textit{Editorial note.} The counted span cites one same-chapter lemma (lemma-constant-exact, already carded in this lane), recorded in assembly.external\_references with canonical tags as navigation pointers.\par\medskip
\noindent\textit{Allowed standard prerequisites (editorial).}\begin{itemize}
\item Exactness of the constant sheaf functor on exact sequences. Stacks Project, Modules on Sites, Lemma lemma-constant-exact (tag 01AC-series; carded locally as S11-sites-mod-constant-exact). Supplies that applying \textbackslash{}underline\{\textbackslash{} \} to the tensored exact sequence yields an exact sequence of section groups.
\item Right exactness of tensor product. Stacks Project, Algebra, section on tensor products. Underlies the step producing M\textasciicircum{}\{\textbackslash{}oplus m\} \textbackslash{}to M\textasciicircum{}\{\textbackslash{}oplus n\} \textbackslash{}to M \textbackslash{}otimes\_\textbackslash{}Lambda Q \textbackslash{}to 0 from the presentation of Q.
\end{itemize}\medskip
\ifdefined\hypertarget\hypertarget{sourcepointer-1}{}\fi
\noindent\textit{Source reference.} Original lemma 093J: sources/stacks-chow/tags.txt, source lines 11517--11517. Cross-chapter standard prerequisite cited by the author inside the counted span; it is used as a pointer and its own proof is not part of this item.\par\medskip
\begin{lemma}
\label{lemma-tensor-with-finitely-presented}
Let $\mathcal{C}$ be a site. Let $\Lambda$ be a ring and let $M$
and $Q$ be $\Lambda$-modules. If $Q$ is a finitely presented
$\Lambda$-module, then we have
$\underline{M \otimes_\Lambda Q}(U) = \underline{M}(U) \otimes_\Lambda Q$
for all $U \in \Ob(\mathcal{C})$.
\end{lemma}

\begin{proof}
Choose a presentation $\Lambda^{\oplus m} \to \Lambda^{\oplus n} \to Q \to 0$.
This gives an exact sequence
$M^{\oplus m} \to M^{\oplus n} \to M \otimes Q \to 0$.
By Lemma \ref{lemma-constant-exact} we obtain an exact sequence
$$
\underline{M}(U)^{\oplus m} \to
\underline{M}(U)^{\oplus n} \to
\underline{M \otimes Q}(U) \to 0
$$
which proves the lemma. (Note that taking sections over $U$ always
commutes with finite direct sums, but not arbitrary direct sums.)
\end{proof}

\end{document}
